% ============================================================================
% INFORME TECNICO - PRACTICA EXPERIMENTAL UNIDAD IV
% Equipo D - Aplicaciones Web
% ============================================================================
\documentclass[12pt,a4paper]{article}

\usepackage[utf8]{inputenc}
\usepackage[T1]{fontenc}
\usepackage[spanish,es-tabla]{babel}
\usepackage{lmodern}
\usepackage[margin=2.5cm]{geometry}
\usepackage{graphicx}
\usepackage{booktabs}
\usepackage{tabularx}
\usepackage{longtable}
\usepackage{array}
\usepackage{xcolor}
\usepackage{float}
\usepackage{caption}
\usepackage{enumitem}
\usepackage{microtype}
\usepackage{listings}
\usepackage{textcomp}
\usepackage[hidelinks]{hyperref}
\hypersetup{colorlinks=true, urlcolor=blue, linkcolor=blue, citecolor=blue}
\hypersetup{pdftitle={Practica Experimental Unidad IV -- Aplicaciones Web -- UTEQ}, pdfauthor={Equipo D -- UTEQ}}
\urlstyle{same}

\setlength{\parskip}{4pt}
\setlength{\parindent}{1cm}

% ---------------------------------------------------------------------------
% Estilo para bloques de codigo (listings): recuadro con fondo, salto de linea
% automatico y separacion del texto circundante
% ---------------------------------------------------------------------------
\definecolor{codebg}{RGB}{248,248,248}
\definecolor{codeframe}{RGB}{170,170,170}
\definecolor{codecomment}{RGB}{120,120,120}

\lstdefinestyle{pfc-code}{%
  basicstyle=\ttfamily\footnotesize,
  columns=fullflexible,
  keepspaces=true,
  showstringspaces=false,
  tabsize=2,
  breaklines=true,
  breakatwhitespace=true,
  postbreak=\mbox{\textcolor{codeframe}{$\hookrightarrow$}\space},
  frame=single,
  framerule=0.4pt,
  rulecolor=\color{codeframe},
  backgroundcolor=\color{codebg},
  xleftmargin=6pt,
  xrightmargin=6pt,
  aboveskip=10pt,
  belowskip=10pt,
  captionpos=b,
  upquote=true
}
\lstset{style=pfc-code}

% ---------------------------------------------------------------------------
% Colores y badges de la caratula (Equipo D)
% ---------------------------------------------------------------------------
\definecolor{verdeuteq}{RGB}{0,92,47}
\definecolor{verdelight}{RGB}{0,128,64}
\definecolor{javacolor}{RGB}{200,95,20}
\definecolor{angularcolor}{RGB}{195,0,47}
\definecolor{postgrescolor}{RGB}{51,103,145}
\definecolor{rediscolor}{RGB}{164,31,27}
\definecolor{nginxcolor}{RGB}{0,150,57}

\newcommand{\badgeJAVA}{\colorbox{javacolor}{\color{white}\small\bfseries\ Java 21 / Spring Boot 3.5.x\ }}
\newcommand{\badgeAngular}{\colorbox{angularcolor}{\color{white}\small\bfseries\ Angular 17+\ }}
\newcommand{\badgePostgres}{\colorbox{postgrescolor}{\color{white}\small\bfseries\ PostgreSQL 16\ }}
\newcommand{\badgeRedis}{\colorbox{rediscolor}{\color{white}\small\bfseries\ Redis Cache/Blacklist\ }}
\newcommand{\badgeNginx}{\colorbox{nginxcolor}{\color{white}\small\bfseries\ Docker / Nginx\ }}

% Marca visible para partes que faltan por completar (se quita al cerrar el documento)
\newcommand{\pendiente}[1]{%
  \par\vspace{4pt}\noindent\textbf{\textcolor{red}{[PENDIENTE: #1]}}\par\vspace{4pt}}

% Figura segura: si falta el archivo de imagen, muestra un recuadro y no rompe la compilacion
\newcommand{\captura}[4]{%
  \begin{figure}[H]
    \centering
    \IfFileExists{#1}{\includegraphics[width=#4]{#1}}{\fbox{\parbox{0.8\textwidth}{\centering Imagen pendiente de subir a Overleaf: \texttt{#1}}}}
    \caption{#2}
    \label{#3}
  \end{figure}}

\begin{document}
% babel[spanish] activa "<" y ">" como shorthand de comillas guillemet («»),
% lo que rompe genericos tipo Observable<ApiResponse<T>> dentro de \texttt.
\shorthandoff{<>}

% ---------------------------------------------------------------------------
% PORTADA (Equipo D, con badges de tecnologia)
% ---------------------------------------------------------------------------
\begin{titlepage}
\thispagestyle{empty}
\centering
\vspace*{0.5cm}
{\large\bfseries Universidad Técnica Estatal de Quevedo (UTEQ)}\\[0.2cm]
{\normalsize Facultad de Ciencias de la Computación}\\[0.1cm]
{\normalsize Carrera de Ingeniería de Software (Rediseño)}\\[0.1cm]
{\normalsize Asignatura: Aplicaciones Web --- Período académico 2026--2027}\\[1cm]
\rule{0.9\linewidth}{1pt}\\[0.4cm]
{\Large\bfseries\color{verdeuteq} Práctica Experimental --- Unidad IV}\\[0.2cm]
{\normalsize\bfseries Modelo Vista Controlador y Servicios Web: API REST, JWT y Consumo de APIs Externas}\\[0.3cm]
\rule{0.9\linewidth}{1pt}\\[0.6cm]
\begin{center}
\badgeJAVA\quad\badgeAngular\\[0.2cm]
\badgePostgres\quad\badgeRedis\quad\badgeNginx\\[0.2cm]
\end{center}
\vspace{0.6cm}
{\normalsize\bfseries Integrantes (Equipo D):}\\[0.1cm]
Castro Espinoza Kevin Moisés\\
Escudero Plaza María del Rosario\\
Loor Medranda Marlon Taylor\\[0.6cm]
{\normalsize\bfseries Docente:}\\
Ing. Gleiston Cicerón Guerrero Ulloa\\[0.6cm]
{\normalsize\bfseries Fecha de entrega:}\\
Domingo, 9 de agosto de 2026\\[0.8cm]
\textbf{Último Commit: Lunes 10 de agosto a las 23:55}\\
{\normalsize\bfseries Repositorio del Proyecto (Rama: \texttt{main}):}\\[0.4cm]
{\small
\textbf{Enlace al repositorio oficial en GitHub:}\\
\url{https://github.com/mloorm14/practica-experimental-unidad-iv}
}\\[0.5cm]
{\small SOFT-A --- 5to Nivel --- Práctica Experimental Unidad IV | Grupal}
\end{titlepage}

\tableofcontents
\newpage

% ---------------------------------------------------------------------------
% RESUMEN
% ---------------------------------------------------------------------------
\section*{Resumen}
\addcontentsline{toc}{section}{Resumen}

El presente trabajo describe la práctica experimental de la Unidad IV de Aplicaciones
Web, desarrollada sobre un sistema de gestión bibliotecaria compuesto por un backend
Spring Boot 3.5.16 (Java 21), un frontend Angular, PostgreSQL 16, Redis y una
infraestructura de despliegue basada en Docker Compose y Nginx. Durante esta unidad el
equipo implementó la versión final de la API REST bajo el prefijo \texttt{/api/v1}, un
formato de respuesta uniforme (\texttt{ApiResponse}), el manejo de errores estandarizado
con \texttt{ProblemDetail}, la autenticación con JWT firmado en HS384 almacenado en una
cookie \texttt{HttpOnly}, la autorización por roles con \texttt{@PreAuthorize}, la
integración con la API externa Open Library y un servicio SOAP contract-first para la
consulta de libros por ISBN. Como verificación, el repositorio cuenta con 48 pruebas
automáticas (unitarias, de repositorio y 12 nuevas de integración HTTP), una prueba de
carga con Apache Bench que alcanzó 986.8 requests por segundo con cache caliente y una
auditoría de seguridad sobre OWASP Top 10:2021 aplicada contra el código real.

\textbf{Palabras clave:} API REST, JWT, Spring Boot, Testcontainers, SOAP, OpenAPI,
seguridad, Docker.

\section*{Abstract}
\addcontentsline{toc}{section}{Abstract}

This work describes the experimental practice of Unit IV of the Web Applications
course, developed on a library management system made of a Spring Boot 3.5.16 backend
(Java 21), an Angular frontend, PostgreSQL 16, Redis and a deployment infrastructure
based on Docker Compose and Nginx. During this unit the team implemented the final
version of the REST API under the \texttt{/api/v1} prefix, a uniform response format
(\texttt{ApiResponse}), standardized error handling with \texttt{ProblemDetail}, JWT
authentication signed in HS384 stored in an \texttt{HttpOnly} cookie, role-based
authorization with \texttt{@PreAuthorize}, the integration with the external Open
Library API and a contract-first SOAP service to look up books by ISBN. As
verification, the repository has 48 automated tests (unit, repository and 12 new HTTP
integration tests), a load test with Apache Bench that reached 986.8 requests per
second with a warm cache and an OWASP Top 10:2021 security audit applied against the
real code.

\textbf{Keywords:} REST API, JWT, Spring Boot, Testcontainers, SOAP, OpenAPI,
security, Docker.

\newpage

% ---------------------------------------------------------------------------
% 1. INTRODUCCION
% ---------------------------------------------------------------------------
\section{Introducción}

La práctica experimental de la Unidad IV continúa el sistema de gestión bibliotecaria
iniciado en la Unidad III. El objetivo general del curso es que el equipo construya una
aplicación web completa con arquitectura MVC, que exponga e integre servicios web REST
y SOAP, y que documente su funcionamiento mediante diagramas de arquitectura y pruebas
verificables.

En esta unidad el trabajo se centró en consolidar el backend y cerrar los objetivos
específicos pendientes. El backend quedó versionado bajo \texttt{/api/v1}, con un
envelope de respuesta uniforme para las respuestas exitosas (\texttt{ApiResponse}) y
\texttt{ProblemDetail} (RFC 7807) para los errores. Se completaron cinco recursos CRUD
(Libro, Autor, Editorial, Idioma y EstadoLibro), la autenticación JWT en cookie
\texttt{HttpOnly} con revocación de tokens mediante una blacklist de \texttt{jti} en
Redis, la autorización por rol con \texttt{@PreAuthorize}, el consumo de la API externa
Open Library con cache-aside y un servicio SOAP contract-first para consultar un libro
por ISBN. El frontend Angular quedó conectado a la API, con login, listado paginado y
CRUD de libros, y la aplicación se despliega completa con Docker Compose (cinco
servicios) detrás de un reverse proxy Nginx.

Este documento presenta el desarrollo realizado, las pruebas y resultados obtenidos, el
marco teórico que sustenta las decisiones técnicas y la evidencia recopilada (prueba de
carga, auditoría de seguridad y pruebas de integración). Al final se incluyen las
reflexiones individuales del equipo, las conclusiones y el trabajo futuro.

\newpage

% ---------------------------------------------------------------------------
% 2. DESARROLLO DEL PROYECTO
% ---------------------------------------------------------------------------
\section{Desarrollo del proyecto}

El repositorio (\url{https://github.com/mloorm14/practica-experimental-unidad-iv})
contiene todo el trabajo de las tres unidades. La Unidad IV se desarrolló sobre la rama
\texttt{main}, con commits convencionales en español (formato
\texttt{feat:/fix:/docs:}), siguiendo la guía interna del equipo que exige el
\texttt{./mvnw -B clean verify} en verde antes de cada commit del backend.

\subsection{Backend}

El backend es una aplicación Spring Boot 3.5.16 sobre Java 21 (\hyperlink{ref:springboot}{Spring Boot, s.f.}), con las siguientes
características implementadas y verificadas en el repositorio:

\begin{itemize}[leftmargin=1.2cm]
  \item \textbf{API versionada}: todos los endpoints viven bajo \texttt{/api/v1}.
        (commit \texttt{e46b0b7}).
  \item \textbf{Envelope uniforme}: las respuestas exitosas usan
        \texttt{ApiResponse<T>} con los campos \texttt{success}, \texttt{data},
        \texttt{message}, \texttt{errors} y \texttt{meta}. Los errores usan
        \texttt{ProblemDetail}, gestionados de forma centralizada en
        \texttt{GlobalExceptionHandler}.
  \item \textbf{Autenticación JWT}: el token se firma con HMAC-SHA384 (\texttt{jjwt}),
        viaja en una cookie \texttt{HttpOnly} con \texttt{SameSite=Strict} y también se
        acepta por header \texttt{Authorization: Bearer} (filtro
        \texttt{JwtAuthenticationFilter}).
  \item \textbf{Revocación de sesión}: al hacer logout, el \texttt{jti} del token se
        guarda en una blacklist de Redis con TTL igual a la vida restante del token
        (\texttt{TokenBlacklistService}).
  \item \textbf{Autenticación por rol}: las operaciones de escritura de los cinco
        recursos usan \texttt{@PreAuthorize("hasRole('ADMIN')")}, con
        \texttt{@EnableMethodSecurity} (\hyperlink{ref:springsecurity}{Spring Security, s.f.}).
  \item \textbf{Cinco recursos CRUD}: Libro (con relación N:M hacia Autor), Autor,
        Editorial, Idioma y EstadoLibro. La gestión del esquema se hace solo con
        migraciones Flyway versionadas (\texttt{V1} a \texttt{V8}) y
        \texttt{hibernate.ddl-auto: validate}.
  \item \textbf{Integración externa}: \texttt{OpenLibraryClient} consulta la API de
        Open Library por ISBN, con cache-aside de 24 horas en un namespace Redis
        separado y un endpoint \texttt{GET /api/v1/libros/\{id\}/enriquecido} que nunca
        falla si el servicio externo no responde.
  \item \textbf{Servicio SOAP}: servicio contract-first con Spring-WS
        (\texttt{libro-catalogo.xsd} $\rightarrow$ JAXB), expuesto en
        \texttt{/ws/libro-catalogo.wsdl}, con manejo de \texttt{SOAP Fault} para ISBN
        inexistente.
  \item \textbf{Documentación}: OpenAPI/Swagger UI en \texttt{/api/documentation}
        (con alias \texttt{/api/docs}), configurado con springdoc-openapi 2.9.0.
\end{itemize}

\subsection{Frontend}

El frontend es una aplicación Angular 17+ conectada a la API mediante el servicio
\texttt{libro.service.ts}, que usa \texttt{Observable<ApiResponse<T>>} y desenvuelve el
\texttt{.data} del envelope antes de llegar a los componentes. Incluye login, listado
paginado de libros con ruta protegida, formulario para crear y editar libros,
eliminación desde el listado y un interceptor de errores 401 con guard de autenticación.

En esta unidad Kevin amplió el frontend con cuatro entregas verificables en el
historial de la rama \texttt{feature/frontend-kevin} (21 commits, mergeados a
\texttt{main} en el commit \texttt{546380e}):

\begin{enumerate}[leftmargin=1.2cm]
  \item \textbf{CRUD completo de Autor} (\texttt{autor.service.ts}, commit
        \texttt{cb8132b}), replicando literalmente el patrón de
        \texttt{libro.service.ts}: cada método tipado como
        \texttt{Observable<ApiResponse<T>>} y desenvuelto con
        \texttt{.pipe(map(r => r.data))}, con las páginas de listado y formulario
        (commits \texttt{8aac2fe} y \texttt{0060bab}).
  \item \textbf{Página de detalle de libro con Open Library} (\texttt{libro-detalle/},
        commit \texttt{ade913d}), que consume
        \texttt{GET /api/v1/libros/\{id\}/enriquecido} y muestra
        \texttt{tituloOpenLibrary}, \texttt{coverUrl}, \texttt{numeroPaginas} y
        \texttt{descripcionOpenLibrary} con manejo explícito de campos \texttt{null}
        (commits \texttt{64772a4} y \texttt{42fd57f} para el modelo y el método del
        servicio).
  \item \textbf{Vistas condicionadas por rol}: el helper \texttt{esAdmin()} del
        \texttt{AuthService} (commit \texttt{5173124}) oculta crear/editar/eliminar
        para el rol \texttt{USER} en libros y autores (commit \texttt{e3af62c}).
  \item \textbf{Progressive Web App}: service worker de Angular
        (\texttt{@angular/service-worker}) y manifest instalable (commits
        \texttt{82779c9}, \texttt{1532b0b} y \texttt{0050774}).
\end{enumerate}

El detalle técnico de los tropiezos reales durante esta implementación (un error de
compilación \texttt{TS1205} al reutilizar un tipo entre servicios, una suposición
incorrecta sobre el código de respuesta del logout) está documentado en la reflexión
individual de Kevin (sección \ref{sec:reflexion-kevin}).

\subsection{Infraestructura}

El despliegue se definió en \texttt{docker-compose.yml} con cinco servicios:
\texttt{postgres} (PostgreSQL 16; \hyperlink{ref:postgres}{PostgreSQL Global Development Group, s.f.}), \texttt{redis} (Redis 7), \texttt{app} (backend
Spring Boot), \texttt{frontend} (Angular servido por Nginx) y \texttt{nginx} (reverse
proxy de entrada, único puerto publicado al host). El \texttt{.env} define
\texttt{JWT\_SECRET} y \texttt{COOKIE\_SECURE} (seguro por defecto; solo se relaja en
desarrollo local). Todos los contenedores levantan con un solo comando
\texttt{docker compose up -d --build}.

\newpage

% ---------------------------------------------------------------------------
% 3. PRUEBAS Y RESULTADOS
% ---------------------------------------------------------------------------
\section{Pruebas y resultados}

\subsection{Pruebas automatizadas}
\label{sec:pruebas}

El proyecto pasó de 36 pruebas (unitarias y de repositorio) a 48, agregando en esta
unidad el paquete \texttt{backend/src/test/java/ec/edu/uteq/pfcbackend/integration/},
con 12 pruebas de integración HTTP end-to-end sobre MockMvc + Testcontainers
(PostgreSQL 16 y Redis 7 reales). Las pruebas se organizaron en dos clases:

\begin{itemize}[leftmargin=1.2cm]
  \item \texttt{AuthIntegracionTest} (5 pruebas): login con credenciales válidas
        devuelve 200 y fija la cookie \texttt{access\_token} con \texttt{HttpOnly};
        login con credenciales inválidas devuelve 401; login con usuario inexistente
        devuelve 401; listar libros sin token devuelve 401; y logout invalida el token
        (el mismo token pasa de 200 a 401 tras el logout).
  \item \texttt{LibroIntegracionTest} (7 pruebas): listar libros con rol USER devuelve
        200; crear un libro con rol USER devuelve 403; crear un libro con rol ADMIN
        devuelve 201; obtener un libro inexistente devuelve 404 con \texttt{ProblemDetail};
        crear un libro con body inválido devuelve 400 con el arreglo \texttt{errors}
        poblado; eliminar un libro con rol USER devuelve 403; y una respuesta exitosa
        conserva el envelope \texttt{ApiResponse}.
\end{itemize}

El resultado del build completo es reproducible con el comando
\texttt{./mvnw -B clean verify}: 48 pruebas ejecutadas, 0 fallos y 0 errores, con
compilación, empaquetado y cobertura JaCoCo incluidos.

\captura{figuras/figura-1-build-verify.png}{Resultado del build \texttt{clean verify}: 48 pruebas ejecutadas, 0 fallos, BUILD SUCCESS.}{fig:build}{0.85\textwidth}

\captura{figuras/figura-2-jacoco.png}{Reporte de cobertura JaCoCo generado por el build.}{fig:jacoco}{0.85\textwidth}

\captura{figuras/figura-3-surefire.png}{Reporte de surefire de las pruebas de integración (\texttt{AuthIntegracionTest} y \texttt{LibroIntegracionTest}).}{fig:surefire}{0.85\textwidth}

\captura{figuras/figura-4-commits-maria.png}{Historial de commits de la rama con las pruebas de integración.}{fig:gitlog}{0.85\textwidth}

\subsection{Prueba de carga con Apache Bench}
\label{sec:carga}

Se probó, con la herramienta \texttt{ab} de Apache (\hyperlink{ref:apache}{The Apache Software Foundation,
s.f.}), el endpoint más usado del sistema (\texttt{GET /api/v1/libros}) contra Nginx, con
20 usuarios concurrentes y 500 requests por corrida, usando un token JWT real. Se
ejecutaron dos corridas: una con cache fría (tras \texttt{FLUSHALL} en Redis) y otra con
cache caliente. La tabla siguiente resume los resultados extraídos del reporte real de
la herramienta.

\begin{table}[H]
\centering
\caption{Resumen de la prueba de carga con Apache Bench (500 requests, 20 usuarios concurrentes).}
\label{tab:benchmark}
\begin{tabular}{lccc}
\toprule
Métrica & Cache frío & Cache caliente & Diferencia \\
\midrule
Requests completados & 500 & 500 & --- \\
Requests fallidos & 0 & 0 & --- \\
Throughput (req/s) & 802.66 & 986.80 & +22.9\% \\
Tiempo medio por request & 24.917 ms & 20.268 ms & $-$18.7\% \\
P50 (mediana) & 23 ms & 19 ms & $-$17.4\% \\
P95 & 32 ms & 27 ms & $-$15.6\% \\
P99 & 36 ms & 30 ms & $-$16.7\% \\
Máximo & 37 ms & 34 ms & $-$8.1\% \\
\bottomrule
\end{tabular}
\end{table}

Ambas corridas terminaron con 0 requests fallidos. La corrida con cache caliente fue más
rápida en todos los percentiles y mantuvo mayor throughput, lo que confirma el beneficio
del patrón cache-aside implementado con Redis.

\subsection{Auditoría de seguridad OWASP Top 10:2021}

La auditoría se aplicó contra el código real del repositorio y, donde fue posible, se
ejecutaron ataques reales con \texttt{curl} contra el sistema en vivo. De ocho controles
auditados, seis están completamente mitigados, uno tiene mitigación parcial (A07, con
una nota concreta sobre fuerza bruta) y uno es un gap real no mitigado (rate limiting,
que también afecta a A07). El resumen de la auditoría es el siguiente:

\begin{table}[H]
\centering
\caption{Resumen de la auditoría OWASP Top 10:2021 aplicada al proyecto (vector, contramedida y evidencia de cada control).}
\label{tab:owasp}
\footnotesize
\setlength{\tabcolsep}{4pt}
\renewcommand{\arraystretch}{1.3}
\begin{tabularx}{\textwidth}{lXXl}
\toprule
Vulnerabilidad & Vector de ataque & Contramedida implementada & Evidencia \\
\midrule
A01 Broken Access Control & Usuario \texttt{USER} intenta un write & \texttt{@PreAuthorize("hasRole('ADMIN')")} + \texttt{@EnableMethodSecurity} & curl $\rightarrow$ 403 \\
A02 Cryptographic Failures & Dump de BD / robo de token & BCrypt (costo 10), JWT HS384, cookie \texttt{Secure} por entorno & hash \texttt{\$2a\$10\$...} \\
A03 Injection & SQL en \texttt{titulo} & Repositorio parametrizado (Spring Data JPA) & curl $\rightarrow$ 0 resultados \\
A04 Insecure Design & Fuerza bruta / DoS de aplicación & \textbf{Gap}: sin rate limiting (Bucket4j disponible) & 5 intentos $\rightarrow$ 5\texttimes 401 \\
A05 Security Misconfiguration & Acceso a \texttt{/actuator/env} & Nginx no enruta \texttt{/actuator}; exposición \texttt{health,info} & \texttt{index.html} / 500 genérico \\
A06 Vulnerable Components & CVE en librerías & Gap parcial: \texttt{jjwt} 0.12.6 $\rightarrow$ 0.13.0 & \texttt{versions:display-dependency-updates} \\
A07 Auth Failures & Reuso de JWT robado post-logout & Blacklist de \texttt{jti} en Redis + TTL & misma cookie 200 $\rightarrow$ 401 \\
XSS & \texttt{onerror} en \texttt{titulo} & Sanitización Angular + cookie \texttt{HttpOnly} & 0 usos de \texttt{innerHTML} \\
\bottomrule
\end{tabularx}
\end{table}

Entre los hallazgos relevantes: los intentos de inyección SQL sobre el parámetro
\texttt{titulo} se tratan como texto literal (0 resultados, la tabla queda intacta); el
login con usuario \texttt{USER} contra una operación de escritura devuelve 403; y la
misma cookie reutilizada después del logout pasa de 200 a 401, verificando la blacklist
de \texttt{jti} en Redis.

\subsection{Verificación de los procedimientos de la guía}

La guía pide verificar, en orden, una lista de procedimientos. El estado real de cada uno
es el siguiente:

\begin{enumerate}[label=\textbf{Paso \arabic*.}, itemsep=2pt]
  \item \textbf{Pruebas automatizadas (10+ pasando, 0 fallos)}: el build
        \texttt{./mvnw -B clean verify} ejecuta 48 pruebas con 0 fallos y 0 errores,
        verificable en las capturas 1 a 4 de la sección \ref{sec:pruebas}.
  \item \textbf{Documentación Swagger y colección Postman}: Swagger UI es accesible en
        \texttt{http://localhost/api/documentation}, con todos los endpoints de los cinco
        recursos CRUD documentados con sus esquemas. La colección de Postman exportada en
        \texttt{docs/postman/SGB-API.postman\_collection.json} contiene \textbf{29
        requests} organizadas en \textbf{6 carpetas} (Auth: 2; Libros: 7; Autores: 5;
        Editoriales: 5; Idiomas: 5; Estados: 5), superando el mínimo de 20 que exige la
        guía.
  \item \textbf{API externa visible en la interfaz}: el endpoint
        \texttt{GET /api/v1/libros/\{id\}/enriquecido} enriquece el libro con los datos
        de Open Library. La respuesta de la API externa se almacena en Redis en el
        namespace \texttt{openlibrary\_isbn} con TTL de 24\,h
        (\texttt{RedisCacheConfig.java}); se verifica con
        \texttt{docker compose exec redis redis-cli KEYS "openlibrary\_isbn*"}, que lista
        las claves cacheadas de cada ISBN consultado, y el cliente distingue
        \texttt{NoEncontrado} de \texttt{ServicioNoDisponible} para ofrecer un mensaje
        amigable en la interfaz cuando el proveedor externo no responde. La
        visualización se implementa en
        \texttt{frontend/src/app/pages/libro-detalle/libro-detalle.html} (componente
        \texttt{LibroDetalle}), que muestra \texttt{tituloOpenLibrary},
        \texttt{coverUrl}, \texttt{numeroPaginas} y \texttt{descripcionOpenLibrary} con
        \texttt{?? 'No disponible'} como valor de respaldo en cada campo, y un
        placeholder ("Portada no disponible.") cuando \texttt{coverUrl} es
        \texttt{null} — de modo que la página nunca se rompe si Open Library no
        responde o el ISBN no existe ahí. Verificado en vivo para ambos casos: el
        libro con ISBN de "1984" (con datos) y "Cien años de soledad" (sin datos),
        ninguno de los dos produjo errores de JavaScript en el navegador.
  \item \textbf{Headers HTTP, carga y despliegue}: la prueba de carga
        (\texttt{ab -n 500 -c 20}) está en la sección \ref{sec:carga}; el despliegue se
        verifica con \texttt{docker compose up -d} (sin errores) y \texttt{docker compose
        ps}, que muestra los cinco servicios \texttt{Up} con sus healthchecks. La
        verificación de los headers de seguridad se ejecutó en vivo contra el sistema
        real (\texttt{curl -I} sobre \texttt{GET /api/v1/libros} con un token JWT
        válido, a través de Nginx):
\begin{lstlisting}[basicstyle=\ttfamily\footnotesize, caption={Headers HTTP reales devueltos por el sistema (\texttt{curl -I}, 11 de agosto de 2026).}]
HTTP/1.1 200
Server: nginx/1.31.3
Date: Tue, 11 Aug 2026 03:29:46 GMT
Content-Type: application/json
Connection: keep-alive
X-Content-Type-Options: nosniff
X-XSS-Protection: 0
Cache-Control: no-cache, no-store, max-age=0, must-revalidate
Pragma: no-cache
Expires: 0
X-Frame-Options: DENY
\end{lstlisting}
        Confirma lo esperado: Spring Security no deshabilita el bloque de headers por
        defecto (\texttt{X-Content-Type-Options: nosniff}, \texttt{X-Frame-Options:
        DENY}, \texttt{Cache-Control} sin cache para respuestas autenticadas).
\end{enumerate}

\newpage

% ---------------------------------------------------------------------------
% 4. TENDENCIAS TECNOLOGICAS
% ---------------------------------------------------------------------------
\section{Tendencias tecnológicas}

\subsection{Jamstack}

Jamstack (\hyperlink{ref:netlify}{Netlify, s.f.}) describe una arquitectura donde el contenido se sirve estático desde una CDN y
la lógica dinámica se consume mediante APIs. Los generadores de sitios estáticos que
materializan esta arquitectura se clasifican por el momento en que generan el HTML:
\textit{build-time} (Next.js estático, Nuxt, Astro) o \textit{runtime} (Next.js con SSR o
ISR), una decisión que afecta directamente a la latencia de la primera carga y al
\textit{time-to-interactive} de la aplicación. Este proyecto no adopta Jamstack completo,
pero comparte una de sus ideas: el frontend Angular se construye y se sirve como
estática a través de Nginx, y toda la dinámica se obtiene de la API REST. La diferencia
es que aquí el reverse proxy y la API conviven en el mismo origen (\texttt{http://localhost}),
lo que evita configuraciones de CORS. El equipo considera Jamstack apropiado para
aplicaciones con poco estado y contenido predecible; para un catálogo con escrituras
frecuentes y roles, el costo de desplegar el frontend y el backend por separado no
aporta una ventaja clara.

\subsection{Progressive Web Apps (PWA)}

Una PWA combina una aplicación web con instalación, funcionamiento offline parcial y
notificaciones (\hyperlink{ref:google}{Google, s.f.}). La guía del SGA pide evaluar esta tendencia y el equipo la
implementó en el frontend: el service worker de Angular
(\texttt{@angular/service-worker}, configurado en \texttt{ngsw-config.json}) y un
manifest instalable (\texttt{manifest.webmanifest}) se verificaron en vivo con el stack
completo dockerizado — el service worker se registra con estado \texttt{activated},
sirviendo \texttt{ngsw-worker.js} a través del mismo reverse proxy Nginx que expone la
API. La auditoría Lighthouse sobre ese mismo stack (\texttt{docs/informe/lighthouse/},
10 de agosto de 2026) obtuvo \textbf{Performance 94}, \textbf{Accessibility 96},
\textbf{Best Practices 100} y \textbf{SEO 82}. El equipo valora la PWA como una mejora
razonable para el frontend actual: el listado de libros puede funcionar offline con un
\texttt{service worker} y la instalación mejora la experiencia en móvil, sin cambiar la
arquitectura del backend.

\subsection{IA generativa}

La IA generativa es la tendencia que más debate genera dentro del curso. El equipo la
usa con una posición crítica: es útil para acelerar tareas concretas (generar
estructuras de prueba, resumir documentación, explorar alternativas de configuración),
pero no reemplaza la verificación. En este proyecto toda decisión de arquitectura se
sostiene con evidencia real del repositorio (ADR, mediciones de Apache Bench, auditoría
OWASP) y todo código generado con apoyo de herramientas se revisa, se ejecuta y se cubre
con pruebas antes de considerarse terminado. El riesgo que el equipo identifica en la
adopción irreflexiva de esta tecnología es que produce texto y código aparentemente
correctos que un humano sin criterio técnico puede incorporar sin verificar, y eso se
nota en informes con referencias inventadas o en código que no compila. La postura del
equipo es usarla como apoyo y mantener la responsabilidad de la revisión en las
personas. En la práctica la herramienta se usó en dos modalidades distintas: los
asistentes de autocompletado (Copilot) para código en contexto dentro del IDE, y el
\textit{scaffolding} de plantillas (por ejemplo, la generación inicial de un controlador
CRUD). Ambas modalidades exponen el mismo límite ético: los modelos no garantizan
originalidad ni trazabilidad de su output, por lo que el equipo verificó el código
generado, lo firmó con su autoría y cubrió cada fragmento con pruebas; no se incorporó
ningún texto ni código sin revisión humana. La responsabilidad final de un error
introducido por una herramienta sigue siendo del equipo, y las decisiones de diseño se
argumentaron con mediciones propias en lugar de citar afirmaciones no verificables.

\newpage

% ---------------------------------------------------------------------------
% 5. MARCO TEORICO
% ---------------------------------------------------------------------------
\section{Marco teórico}

\subsection{MVC y comparativa de frameworks}

El patrón Modelo-Vista-Controlador fue descrito originalmente por Trygve Reenskaug en
1979 dentro del entorno Smalltalk-80: el \textit{modelo} gestiona el estado de la
aplicación, la \textit{vista} es la representación de ese estado y el \textit{controlador}
es el manejador de las entradas del usuario. El modelo web adaptó esta separación: el
controlador ya no recibe eventos de teclado o ratón sino solicitudes HTTP; la vista se
renderiza en el servidor con un motor de plantillas o en el cliente (SPA); y el modelo
pasa a incluir repositorios y servicios, no solo el estado en memoria. De esa adaptación
surgieron variantes como MVP (\textit{Model-View-Presenter}), donde la vista no tiene
referencia directa al modelo, y MVVM (\textit{Model-View-ViewModel}), que usa un enlace de
datos bidireccional y es el patrón de facto en SPA modernas (\hyperlink{ref:fowler}{Fowler, 2002}) \hyperlink{ieee2}{[2]}. El backend de
este proyecto aplica MVC en su forma en capas (controlador, servicio, repositorio), tal
como documenta el ADR-001 de la Unidad III, y las entidades mapeadas con JPA/Hibernate
actúan como el modelo.

Todos los frameworks MVC web usan el patrón \textit{Front Controller}: un punto de
entrada único (en PHP, \texttt{index.php}; en .NET, \texttt{Program.cs}; en Spring Boot,
\texttt{DispatcherServlet}). Centralizar la entrada permite manejar errores,
autenticación, logging y enrutamiento en un solo lugar, en vez de repetirlos en cada
página. En este proyecto el ciclo completo de una petición es el siguiente: (1) la
solicitud entra por Nginx, que enruta \texttt{/api/} hacia el backend; (2)
\texttt{DispatcherServlet} selecciona el controlador a partir de la anotación
\texttt{@RequestMapping}; (3) la cadena de filtros de Spring Security
(\texttt{JwtAuthenticationFilter}) extrae el JWT de la cookie o del header
\texttt{Authorization}, valida la firma y consulta la blacklist en Redis, estableciendo el
contexto de seguridad; (4) el controlador recibe el request y valida el body con
\texttt{@Valid}; (5) el servicio aplica la lógica de negocio y el repositorio de Spring
Data JPA consulta PostgreSQL, con el cache-aside en Redis interviniendo en las lecturas
marcadas con \texttt{@Cacheable}; (6) la respuesta se serializa en el envelope uniforme
\texttt{ApiResponse} y se devuelve con el código HTTP correspondiente. Los pasos 4 a 6 se
materializan en el código real siguiente (el controlador, la capa de servicio con cache y
el envelope):

\begin{lstlisting}[language=Java, caption={Controlador, servicio y cache del listado (fragmento real de \texttt{LibroController.java} y \texttt{LibroServiceImpl.java}).}]
@RestController
@RequestMapping("/api/v1/libros")
public class LibroController {
    @GetMapping
    @Operation(summary = "Listar libros, con filtro opcional por titulo")
    public ApiResponse<List<LibroResponse>> listar(
            @RequestParam(required = false) String titulo, Pageable pageable) {
        Page<LibroResponse> pagina = libroService.listar(titulo, pageable);
        return new ApiResponse<>(pagina.getContent(), pagina);
    }
}

@Service
public class LibroServiceImpl {
    @Cacheable(value = CACHE_LISTADO,
        key = "#titulo + '-' + #pageable.pageNumber + '-' + #pageable.pageSize + ...")
    public Page<LibroResponse> listar(String titulo, Pageable pageable) {
        return (titulo == null || titulo.isBlank())
            ? libroRepository.findAll(pageable)
            : libroRepository.findByTituloContainingIgnoreCase(titulo, pageable);
    }
}
\end{lstlisting}

Sobre la elección del framework, el ADR-005 compara las tres opciones que el propio PDF de
la guía menciona como válidas para este curso. La tabla siguiente extiende esa comparación
a los cinco frameworks que la guía exige considerar (Laravel, Symfony, Spring Boot,
ASP.NET Core y Django), con los criterios solicitados; la fila de rendimiento usa la
misma fuente que el ADR-005 (TechEmpower, Round 23, prueba Fortunes), por ser la ronda
publicada más reciente antes del cierre del proyecto, y el resto de los criterios se
tomaron de la documentación oficial de cada framework (\hyperlink{ref:otwell}{Otwell, 2024}) \hyperlink{ieee1}{[1]}:

\begin{table}[H]
\centering
\caption{Comparativa de frameworks MVC (criterios de la guía, datos de ADR-005 y documentación oficial).}
\label{tab:comparativa}
\footnotesize
\setlength{\tabcolsep}{4pt}
\renewcommand{\arraystretch}{1.3}
\begin{tabularx}{\textwidth}{lXXXXX}
\toprule
Criterio & Laravel 11.x & Symfony 7.x & Spring Boot 3.x & ASP.NET Core 8.x & Django 5.x \\
\midrule
Paradigma & Convención sobre config. & Explícito & Explícito & Intermedio & Convención sobre config. \\
ORM incluido & Eloquent & Doctrine & Spring Data JPA & EF Core & Django ORM \\
Motor de plantillas & Blade & Twig & Thymeleaf & Razor & Django Templates \\
Ecosistema & Paquetes (Packagist) & Paquetes & Starters (Maven) & NuGet & Apps (PyPI) \\
Rendimiento (TechEmpower R23, Fortunes) & ~16\,800 req/s & no medido en el ADR & ~243\,639 req/s & ~609\,966 req/s & no medido en el ADR \\
Adopción LATAM & Alta (agencias/freelance) & Media & Alta (empresarial/banca) & Media & Media \\
\bottomrule
\end{tabularx}
\end{table}

El equipo eligió Spring Boot principalmente por tres razones documentadas: la experiencia
previa real con Java, el tipado estático (que detecta errores en tiempo de compilación,
como el desajuste de tipos detectado en el ADR-004) y un ecosistema de testing maduro
(Testcontainers levantando un PostgreSQL real en las pruebas, sin mocks de base de
datos). Se reconoce en el propio ADR-005 la contraparte honesta: Java es más verboso que
Laravel para el mismo CRUD (cada recurso requirió siete archivos), y la configuración
inicial de Spring es más pesada que la de un microframework. La elección se justifica
además por el rendimiento medido en este mismo proyecto (P99 de 30 a 36 ms en la prueba de
carga), que confirma que el rendimiento bruto de ASP.NET Core no era la restricción activa
del PFC (ADR-005).

\subsection{REST, JWT y OpenAPI}

HTTP es un protocolo sin estado por diseño: cada petición es independiente y el servidor
no guarda memoria de interacciones previas (\hyperlink{ref:fielding2000}{Fielding, 2000}; \hyperlink{ref:fielding2022}{Fielding et al., 2022}) \hyperlink{ieee3}{[3]}. REST
se apoya en seis principios que la API de este proyecto aplica (o relega, cuando no
aplica) de forma explícita y verificable:

\begin{itemize}
  \item \textbf{Cliente-Servidor}: separación total entre el frontend Angular y la API;
        el backend nunca genera HTML y el frontend nunca accede a la base de datos.
  \item \textbf{Stateless}: el servidor no guarda sesión; la identidad viaja en el JWT de
        cada petición. Esto es lo que habilita escalar horizontalmente sin
        \textit{sticky sessions}.
  \item \textbf{Cacheable}: cada recurso declara con \texttt{Cache-Control} si su
        respuesta puede almacenarse; el cache-aside en Redis (un TTL por dominio) es la
        implementación concreta de este principio.
  \item \textbf{Interface uniforme}: identificación de recursos por URI, representación
        JSON autodescriptiva (\texttt{ApiResponse}), mensajes autodescriptivos y HATEOAS
        como nivel de madurez que el proyecto alcanza parcialmente (enlaces en el
        envelope) sin hipermedia completa.
  \item \textbf{Layered System}: frontend, Nginx, backend, Redis y PostgreSQL forman una
        pila en la que cada capa solo conoce a su vecina; Nginx esconde la topología
        interna al cliente.
  \item \textbf{Code-on-Demand}: opcional por definición; no se usa, porque introducir
        código descargado complica la seguridad sin aportar al PFC.
\end{itemize}

La jerarquía de la URI es el contrato de esa interface uniforme, y su diseño sigue las
reglas de la arquitectura REST (\hyperlink{ref:richardson}{Richardson \& Ruby, 2007}) \hyperlink{ieee4}{[4]}: sustantivos en plural
(\texttt{/libros}, no \texttt{/getLibros}), anidamiento solo para subrecursos legítimos,
sin verbos en la URL y con el versionado en el prefijo \texttt{/api/v1}. El diseño
completo de la API es:

\begin{itemize}
  \item \texttt{POST /api/v1/auth/login} y \texttt{POST /api/v1/auth/logout}
  \item \texttt{GET /api/v1/libros} (con filtros), \texttt{GET /api/v1/libros/\{id\}},
        \texttt{POST /api/v1/libros}, \texttt{PUT /api/v1/libros/\{id\}},
        \texttt{DELETE /api/v1/libros/\{id\}} y
        \texttt{GET /api/v1/libros/\{id\}/enriquecido} (consume la API externa)
  \item \texttt{GET|POST|PUT|DELETE /api/v1/autores/\{id\}}, lo mismo para
        \texttt{/editoriales}, \texttt{/idiomas} y \texttt{/estados}
  \item \texttt{GET /actuator/health} para el healthcheck de Docker
\end{itemize}

La estructura del JWT es la de la especificación RFC~7519 (\hyperlink{ref:jones2015}{Jones et al., 2015}): un
\textbf{Header} con el algoritmo (\texttt{HMAC-SHA384} en este proyecto),
un \textbf{Payload} con los claims \texttt{sub} (id de usuario), \texttt{username},
\texttt{roles}, \texttt{iat} y \texttt{exp}, y una \textbf{Signature} calculada con el
secreto de 48 bytes. El flujo es: \texttt{login} valida credenciales contra PostgreSQL →
el servidor firma el JWT y lo entrega en una cookie \texttt{HttpOnly} con
\texttt{SameSite=Strict} (también aceptado por \texttt{Authorization: Bearer}) → en cada
request autenticada, \texttt{JwtAuthenticationFilter} valida la firma y la expiración y
establece el \texttt{SecurityContext} → en el \texttt{logout}, el \texttt{jti} entra en
una blacklist de Redis (\hyperlink{ref:redis}{Redis, s.f.}) con TTL. La cookie \texttt{HttpOnly} no es legible
desde JavaScript, lo que elimina el vector de robo por XSS que existiría con
\texttt{localStorage}; y la blacklist resuelve el problema estructural del JWT de que un
token firmado no puede revocarse (la estrategia de corto TTL + revocación explícita se
documenta en la guía como la vía correcta).

OpenAPI (\hyperlink{ref:openapi}{OpenAPI Initiative, 2023}) actúa como el contrato entre frontend y backend: la
especificación se genera con springdoc-openapi a partir de las anotaciones de los
controladores y se sirve en \texttt{/api/documentation}. Para cada endpoint la
especificación declara \texttt{operationId}, parámetros de ruta y query, el
\texttt{requestBody} con su \texttt{reference} al esquema y las respuestas por código con
sus esquemas, de modo que Angular puede generar su cliente tipado desde el mismo
documento. Para visualizar ese contrato se compararon dos alternativas: \textbf{Swagger
UI}, interactivo (permite ejecutar las peticiones desde el navegador), y \textbf{Redoc},
estático y pensado como documentación de lectura; el proyecto usa Swagger UI por su
interactividad y lo expone en \texttt{/swagger-ui.html}. La documentación queda
sincronizada con el código, a diferencia de documentación escrita a mano que tiende a
quedar desactualizada. El fragmento siguiente es la entrada que springdoc-openapi genera
para el login (resumida), y muestra los elementos que la guía pide identificar:
\texttt{operationId}, parámetros, \texttt{requestBody} con su esquema y las respuestas
por código:

\begin{lstlisting}[caption={Documentación OpenAPI del endpoint de login (generada por springdoc-openapi, resumida).}]
paths:
  /api/v1/auth/login:
    post:
      operationId: login
      summary: Autentica un usuario y devuelve un JWT
      requestBody:
        required: true
        content:
          application/json:
            schema:
              $ref: '#/components/schemas/LoginRequest'   # username, password
      responses:
        '200':
          description: Login exitoso; el JWT viaja en la cookie access_token
          headers:
            Set-Cookie:
              schema: { type: string }
        '401':
          description: Credenciales invalidas
          content:
            application/json:
              schema:
                $ref: '#/components/schemas/ProblemDetail'
components:
  securitySchemes:
    bearerAuth:
      type: http
      scheme: bearer
      bearerFormat: JWT
\end{lstlisting}

\subsection{SOAP vs.\ REST}

SOAP es un protocolo de mensajería basado en XML con un contrato formal (WSDL)
(\hyperlink{ref:w3c}{World Wide Web Consortium, 2000}). Estructuralmente, un mensaje SOAP se compone de un
\texttt{Envelope} que envuelve al \texttt{Header} (metadatos, seguridades WS-Security) y
al \texttt{Body} (la operación), y los errores viajan como \texttt{SOAP Fault}; cada
request indica la operación a ejecutar con el header \texttt{SOAPAction}, y el transporte
no se limita a HTTP, sino que puede ser SMTP o TCP. REST, en cambio, es un estilo
arquitectónico sobre HTTP en el que los recursos se manipulan con los verbos del protocolo
(\hyperlink{ref:fielding2000}{Fielding, 2000}) \hyperlink{ieee3}{[3]}. La comparación entre frameworks y protocolos que aquí se resume sigue
la taxonomía de referencia del SWEBOK en su edición 4.0 (\hyperlink{ref:washizaki}{Washizaki et al., 2024}) \hyperlink{ieee5}{[5]}, que
estructura el cuerpo de conocimiento en ingeniería de software; el proyecto implementa
ambos protocolos, lo que permite compararlos con evidencia real.

\begin{table}[H]
\centering
\caption{Comparación SOAP vs.\ REST aplicada a los servicios implementados.}
\label{tab:soap-rest}
\footnotesize
\setlength{\tabcolsep}{4pt}
\renewcommand{\arraystretch}{1.3}
\begin{tabularx}{\textwidth}{lXX}
\toprule
Criterio & SOAP (implementado) & REST (implementado) \\
\midrule
Formato de mensaje & XML (SOAP Envelope) & JSON \\
Tipado & Estricto (XSD, contract-first) & Definido por el esquema OpenAPI \\
Overhead de serialización & Alto (namespaces) & Bajo \\
Contrato & WSDL (contract-first) & OpenAPI (descripción) \\
Descubrimiento & WSDL auto-generado & Swagger UI \\
Transporte & HTTP (también SMTP/JMS) & HTTP(S) \\
Manejo de errores & SOAP Fault (HTTP 500) & Códigos de estado + ProblemDetail \\
Cacheabilidad & Difícil (POST) & Explícita (GET cacheable) \\
Tamaño de payload (medido) & 564 bytes / 8 campos & 550 bytes / 16 campos \\
Seguridad & WS-Security & HTTPS + JWT \\
Complejidad de implementación & Alta & Baja-media \\
Interoperabilidad & Alta (sistemas corporativos) & Amplia en la web \\
Idempotencia & No definida por SOAP & GET/PUT/DELETE por diseño \\
Facilidad desde JavaScript & Baja (XML + WSDL) & Alta (JSON nativo) \\
Casos de uso actuales (2025) & Banca, sector público, SRI & APIs web, móvil, SPA \\
\bottomrule
\end{tabularx}
\end{table}

La medición real de payload del mismo caso de uso (consultar el libro ``1984'') arroja
un resultado interesante: la respuesta REST es 14 bytes más pequeña que la SOAP a pesar
de traer el doble de campos, porque cada elemento XML necesita etiquetas de apertura y
cierre con namespace. La respuesta SOAP para un ISBN inexistente devuelve un
\texttt{SOAP Fault} bien formado con el mensaje exacto de la causa, mientras el
equivalente REST devuelve un \texttt{ProblemDetail} con su código de estado. El equipo
mantiene el servicio SOAP público a propósito, como demostración del resultado de
aprendizaje que exige exponer ambos tipos de servicio, y no como recurso de negocio que
necesite protección por rol.

El SOAP no es una tecnología muerta en 2025: sigue siendo el estándar donde la
interoperabilidad entre instituciones grandes importa más que la simplicidad. Los casos
de uso vigentes se concentran en la banca (servicios de core bancario entre
instituciones, compensación y clearing), en el sector público ecuatoriano (intercambio
de datos entre instituciones del Estado a través de sus \textit{service bus}) y en el
SRI, que publica servicios SOAP/WSDL para la validación de comprobantes electrónicos;
son sistemas donde los clientes están controlados (no se consumen desde el navegador) y
donde el contrato WSDL firmado reduce ambigüedad contractual. Ninguno de esos casos
justifica hoy una API pública nueva en SOAP, pero un desarrollador ecuatoriano los va a
encontrar en su carrera profesional, y por eso esta práctica los implementa.

Para consumir una API externa, la guía pide comparar los clientes HTTP de los tres
stack que se comparan en la Unidad III: \textbf{GuzzleHTTP} (PHP), \textbf{HttpClient}
(.NET) y \textbf{RestTemplate/WebClient} (Java). Los tres cubren peticiones GET/POST con
headers y JSON, pero se diferencian en tres dimensiones prácticas: los \textbf{timeouts}
de conexión y respuesta (configurables en los tres, en este proyecto 5\,s de conexión y
10\,s de respuesta vía \texttt{WebClientConfig}), el \textbf{retry con backoff} ante
fallos transitorios (Guzzle y HttpClient lo traen como middleware nativo; en WebClient
hay que componerlo con Reactor) y la \textbf{concurrencia}: RestTemplate bloquea el hilo
mientras espera la respuesta, mientras que WebClient es reactivo y no bloquea, lo que
escala mejor bajo carga concurrente. El fragmento siguiente es el consumo real de Open
Library en el endpoint \texttt{/enriquecido}, y muestra el manejo explícito de los tres
tipos de fallo que una API externa puede producir:

\begin{lstlisting}[language=Java, caption={Consumo de API externa con manejo de errores (fragmento de \texttt{OpenLibraryClient.java}).}]
@Cacheable(value = "openlibrary_isbn", key = "#isbn",
        unless = "#result.getClass().getSimpleName().equals('ServicioNoDisponible')")
public OpenLibraryResult buscarPorIsbn(String isbn) {
    String isbnNormalizado = isbn.replaceAll("[^0-9Xx]", "");
    try {
        OpenLibraryResponse datos = openLibraryWebClient.get()
                .uri("/isbn/{isbn}.json", isbnNormalizado)
                .retrieve()
                .bodyToMono(OpenLibraryResponse.class)
                .timeout(TIMEOUT_RESPUESTA)   // 10 s
                .block();
        return datos != null ? new Encontrado(datos) : new NoEncontrado();
    } catch (WebClientResponseException.NotFound ex) {
        return new NoEncontrado();              // 404: dato definitivo, se cachea
    } catch (WebClientResponseException ex) {
        return new ServicioNoDisponible();      // 5xx del proveedor
    } catch (WebClientRequestException ex) {
        return new ServicioNoDisponible();      // red caida, DNS, conexion
    } catch (RuntimeException ex) {
        if (ex.getCause() instanceof TimeoutException) {
            return new ServicioNoDisponible();  // no responde a tiempo
        }
        throw ex;
    }
}
\end{lstlisting}

La distinción entre \texttt{NoEncontrado} y \texttt{ServicioNoDisponible} no es un
detalle: el cache-aside con la anotación \texttt{@Cacheable} de Spring Cache (\hyperlink{ref:redis}{Redis, s.f.})
guarda por defecto la primera respuesta de cada ISBN, y si se cacheara un fallo
transitorio, la API quedaría "congelada" en estado de error durante todo el TTL. Por eso
el \texttt{unless} impide cachear \texttt{ServicioNoDisponible} pero sí cachea el 404,
que es un dato definitivo. Cachear respuestas externas está justificado por tres
razones: respetar el \textit{rate limiting} del proveedor, reducir la latencia
percibida (una consulta externa puede tardar cientos de milisegundos) y tolerar la
indisponibilidad del proveedor. El TTL se decide por dominio: datos meteorológicos
(cambian cada hora) con TTL corto de 10 minutos, tipos de cambio (1 hora) y datos
prácticamente estáticos como el metadato de un ISBN (24 horas); un TTL único para todos
los dominios sería incorrecto para cualquiera de ellos.

\subsection{Seguridad, rendimiento y despliegue}

El proyecto aplica buenas prácticas de seguridad verificadas en la auditoría OWASP
(\hyperlink{ref:owasp}{OWASP Foundation, 2021}) \hyperlink{ieee6}{[6]}. Las contraseñas se almacenan con BCrypt (costo 10, hash
\texttt{\$2a\$10\$...}); el JWT se firma con HMAC-SHA384 (secreto de 48 bytes base64),
rechazando tokens sin firma o con algoritmo distinto; la cookie de sesión es segura por
defecto y condicionada por entorno (antes de esta unidad estaba fijada en
\texttt{false}); y las consultas se ejecutan parametrizadas por Spring Data, sin
concatenación de strings SQL. El diseño defiende la capa HTTP con los headers que Spring
Security habilita por defecto (el bloque de headers no está deshabilitado en
\texttt{SecurityConfig}): \texttt{X-Content-Type-Options: nosniff},
\texttt{X-Frame-Options: DENY} y \texttt{Cache-Control} en las respuestas, verificados en
vivo con \texttt{curl -I} en la sección \ref{sec:pruebas} (Verificación de los
procedimientos de la guía, paso 4).

Sobre rendimiento, la prueba de carga siguió la metodología de la guía (concurrencia
\texttt{-c}, volumen \texttt{-n} y endpoint objetivo): el comando exacto fue
\texttt{ab -n 500 -c 20 -H "Authorization: Bearer <token>" http://nginx/api/v1/libros},
ejecutado desde un contenedor descartable en la misma red de Docker para no instalar nada
en el host (\hyperlink{ref:apache}{Apache Software Foundation, s.f.}). El endpoint más usado respondió con
\textbf{986.80 req/s con cache caliente} frente a \textbf{802.66 req/s con cache fría}, y
todos los percentiles mejoraron: P50 de 23 a 19\,ms, P95 de 32 a 27\,ms, P99 de 36 a
30\,ms. El \textbf{error rate fue 0\% en ambas corridas} (1000 requests totales). Estos
números sustentan que, para el volumen esperado de un sistema académico, la configuración
actual no presenta cuellos de botella. La interpretación de un sistema que degrada
suavemente bajo carga y no produce errores en picos sigue el marco de \textit{Release
It!} (\hyperlink{ref:nygard}{Nygard, 2018}) \hyperlink{ieee7}{[7]}, que separa la resistencia estructural (la aplicación no se cae ante
fallos) del rendimiento puro.

El despliegue se define con Docker Compose (\hyperlink{ref:docker}{Docker, s.f.}): PostgreSQL, Redis, backend, frontend y Nginx
como reverse proxy. Las relaciones entre servicios son explícitas en el compose: la app
depende de \texttt{postgres} y \texttt{redis} (que llegan a estar saludables vía
healthchecks antes de arrancar), y \texttt{nginx} es el único punto de entrada al mundo
exterior. Nginx enruta \texttt{/api/} hacia el backend y \texttt{/} hacia el frontend, y
no expone \texttt{/actuator/*}, lo que impide acceder a endpoints de monitorización desde
el exterior; además, el backend solo habilita \texttt{health} e \texttt{info} de
Actuator. Docker aporta tres beneficios concretos a este proyecto: \textbf{entornos
reproducibles} (el mismo \texttt{docker-compose.yml} levanta el sistema idéntico en
cualquier máquina, sin configurar JDK, Node, PostgreSQL ni Redis a mano),
\textbf{aislamiento de dependencias} (cada servicio encapsula sus librerías y versiones,
evitando el clásico \texttt{it works on my machine}) y la base para \textbf{CI/CD} (los
pipelines pueden construir la misma imagen que se ejecuta en producción). Los secretos y
valores de configuración se inyectan con variables de entorno a través de
\texttt{environment:} en el compose y \texttt{\$\{VARIABLE\}} en \texttt{application.yml}
(por ejemplo \texttt{JWT\_SECRET}, \texttt{COOKIE\_SECURE}, \texttt{DB\_*}); nunca en el
código fuente, porque un archivo con secretos hardcodeados se filtra si el repositorio se
hace público o se comparte, y no puede cambiar entre entornos sin recompilar. El
despliegue completo se levanta con \texttt{docker compose up -d} y se verifica con
\texttt{docker compose ps}, que muestra los cinco servicios \texttt{Up} con sus
healthchecks.

\newpage

% ---------------------------------------------------------------------------
% 6. REFLEXIONES INDIVIDUALES
% ---------------------------------------------------------------------------
\section{Reflexiones individuales}

\subsection{Loor Medranda Marlon Taylor (tech lead / backend / infraestructura)}

A lo largo de este proyecto tuve que corregir tres errores propios que, revisados en
retrospectiva, dejan lecciones más útiles que cualquier cosa que hubiera aprendido si
todo hubiese salido bien a la primera.

El primero fue un login roto en producción local, causado por versionar la API a
\texttt{/api/v1} sin actualizar en el mismo cambio los matchers de
\texttt{SecurityConfig}, que seguían apuntando a las rutas sin versionar (commit
\texttt{8e73f02}). El síntoma no apareció en ningún test unitario ni de compilación:
apareció al probar el flujo real en vivo, porque \texttt{SecurityConfig} y los
controladores viven en archivos distintos y nada obliga a que cambien juntos. La
lección no es "tener cuidado", que es inútil como lección porque no es accionable — es
que un cambio de contrato de API (una ruta, un prefijo, un nombre de campo expuesto)
toca más superficie de la que el diff muestra a primera vista, y que verificar en vivo
después de cada cambio estructural — no solo correr los tests — no es un paso opcional
ni redundante, incluso cuando "solo" se está agregando un prefijo de versión.

El segundo fue un error de comparación \texttt{==} en vez de \texttt{.equals()} dentro
de una expresión SpEL de \texttt{@Cacheable(unless=...)} en \texttt{OpenLibraryClient}
(corregido en el commit \texttt{c80426b}). Esa condición determinaba si un resultado de
Open Library debía cachearse o no, y en las pruebas manuales que hice antes de la
corrección "funcionaba" — porque el caso que probé casualmente caía del lado correcto
de la comparación, no porque la comparación fuera correcta. La lección real ahí es que
pasar una prueba manual no confirma que el mecanismo subyacente sea correcto; confirma
que ese caso particular no expuso el defecto. Un lenguaje/motor de expresiones que no
garantiza comparación por identidad de referencia versus por valor es exactamente el
tipo de detalle que un test "parece que funciona" no detecta, y que solo aparece si se
piensa explícitamente en el contrato del lenguaje, no en el resultado observado una vez.

El tercero lo encontramos nosotros mismos, auditando nuestro propio código con el mismo
rigor que le exigiríamos a un tercero: una cookie de sesión con \texttt{secure(false)}
hardcodeado en \texttt{AuthController}, sin ningún mecanismo que lo condicionara al
entorno de despliegue (corregido en el commit \texttt{aa3e6b3}, ya con
\texttt{app.cookie.secure} configurable y seguro por defecto). Nadie fuera del equipo lo
señaló; lo encontramos porque decidimos hacer la auditoría OWASP contra el código real
en vez de tratarla como un formalismo documental para cumplir un punto del PDF. Esa es
la lección que más me deja este bloque: revisar el propio trabajo con la misma
rigurosidad que se le exigiría a un tercero no es un ejercicio vacío ni redundante —
encuentra fallos reales que de otra forma solo aparecerían en producción, o peor, no
aparecerían nunca porque nadie los busca.

Si empezara este proyecto de nuevo, escribiría desde el día uno un checklist corto de
"qué verificar en vivo tras cada cambio de contrato de API" (autenticación, cada rol
relevante, al menos un caso de error esperado) en vez de descubrir la necesidad de ese
checklist de forma reactiva, después de que un bug de ese tipo ya llegó a producción
local una vez.

\subsection{Escudero Plaza María del Rosario (tests de integración)}

En esta unidad mi tarea fue escribir los tests de integración del backend. La guía pedía
mínimo 10 tests y dejé 12, repartidos en dos clases: una para autenticación (login,
logout y acceso sin token) y otra para el CRUD de libros y los permisos por rol.

Lo que más me costó no fue escribir los tests, sino hacer que funcionaran. Mi
computadora tiene Java 25 instalado y el proyecto usa Java 21, así que al principio los
tests ni siquiera compilaban: Lombok no generaba el código y tuve que correr el build con
una bandera especial (\texttt{-Dmaven.compiler.proc=full}). Ahí entendí que el problema
no era el código, sino que mi entorno no coincidía con el del proyecto.

También asumí cosas que no eran. Por ejemplo, pensé que el logout devolvía un 200 con un
cuerpo, pero en realidad devuelve 204 (sin contenido). Y al probar la creación de
libros, el test fallaba porque el formulario necesita más campos de los que yo le
mandaba (año de publicación, editorial, idioma, estado y stock). Tuve que leer bien el
DTO y el controlador para arreglarlo. Eso me enseñó a revisar el código antes de asumir
cómo funciona un endpoint.

Otro problema fue con los contenedores de prueba. Al separar los tests en dos clases,
compartí los mismos contenedores de base de datos y Redis, y cuando la primera clase
terminaba, la segunda ya no podía conectarse y todo devolvía 500. La solución fue darle
a cada clase sus propios contenedores. Me quedó claro que los tests tienen que estar
aislados entre sí.

Al final el proyecto completo pasa con 48 tests y cero fallos. Si tuviera que repetir la
unidad, empezaría leyendo el código de cada endpoint antes de escribir el primer test.

\subsection{Castro Espinoza Kevin Moisés (frontend)}
\label{sec:reflexion-kevin}

Mi parte de esta unidad fue el frontend Angular: implementar el CRUD de Autor, integrar
Open Library en la vista de detalle, condicionar la UI por rol y convertir la app en
PWA. El primer tropiezo fue técnico y muy concreto: al extraer la interfaz
\texttt{PageResponse} a un modelo compartido para reutilizarla entre
\texttt{libro.service.ts} y el nuevo \texttt{autor.service.ts}, el build falló con
\texttt{TS1205} porque el proyecto usa \texttt{isolatedModules} y TypeScript exige
\texttt{export type} al re-exportar tipos. Fue un recordatorio directo de que "copiar el
patrón" de un archivo existente no basta si no se entiende la configuración del
compilador que lo sostiene.

Aprender a desenvolver el \texttt{ApiResponse} con \texttt{.pipe(map(r => r.data))} en
cada método, tal como ya hacía \texttt{libro.service.ts}, me dejó claro por qué el
contrato del backend (envelope en éxito, \texttt{ProblemDetail} en error) se traduce en
que el componente nunca ve el JSON crudo. También entendí el punto de la guía sobre los
campos de Open Library que pueden venir \texttt{null}: el DTO
\texttt{LibroEnriquecidoResponse} documenta que si el ISBN no existe en Open Library la
respuesta sigue siendo 200, y la vista debía manejar ese caso con placeholders en vez de
romperse.

Finalmente, el rol que viaja en el body del login (\texttt{LoginResponse(username,
rol)}) se guarda en el estado del \texttt{AuthService} y condiciona crear/editar/eliminar
con \texttt{*ngIf}; la cookie \texttt{HttpOnly} hace inviable decodificar el JWT en el
frontend, así que confiar en el body del login era lo correcto. La PWA se agregó
manualmente porque \texttt{ng add @angular/pwa} no resolvió el esquema con el builder de
aplicación (clave \texttt{browser} en \texttt{angular.json}); el reporte Lighthouse
quedó en \texttt{docs/informe/lighthouse} con performance 94 y best practices 100.

\newpage

% ---------------------------------------------------------------------------
% 7. CONCLUSIONES
% ---------------------------------------------------------------------------
\section{Conclusiones}

\begin{enumerate}[leftmargin=1.2cm]
  \item La API REST quedó completa y versionada bajo \texttt{/api/v1}, con un formato de
        respuesta uniforme y errores estandarizados. Esto hace que los clientes (frontend,
        SOAP, pruebas) traten la API de forma consistente.
  \item La autenticación JWT en cookie \texttt{HttpOnly} con blacklist de \texttt{jti}
        en Redis cumple el requisito de sesión sin estado del servidor y permite revocar
        tokens de forma efectiva, verificado en pruebas: el mismo token pasa de 200 a 401
        después del logout.
  \item La autorización por rol (\texttt{@PreAuthorize}) se verifica con pruebas de
        integración: un usuario \texttt{USER} recibe 403 en las operaciones de escritura
        mientras que \texttt{ADMIN} las ejecuta con 201.
  \item El proyecto expone e integra servicios web REST y SOAP. La comparación con datos
        reales muestra que la respuesta REST es más pequeña que la SOAP incluso trayendo
        el doble de campos, y que el servicio SOAP devuelve \texttt{SOAP Fault} bien
        formado para errores.
  \item La prueba de carga y la auditoría OWASP aportan evidencia real de rendimiento y
        seguridad: 986.80 req/s con cache caliente sin requests fallidos, y seis de ocho
        controles OWASP completamente mitigados, con los gaps documentados de forma
        honesta (rate limiting, dependencias desactualizadas).
  \item Las pruebas automatizadas pasaron de 36 a 48 con la incorporación de 12 tests de
        integración HTTP, y todo el build (\texttt{./mvnw -B clean verify}) termina en
        verde, lo que da confianza para continuar evolucionando el sistema.
\end{enumerate}

% ---------------------------------------------------------------------------
% 8. TRABAJO FUTURO
% ---------------------------------------------------------------------------
\section{Trabajo futuro}

\begin{itemize}[leftmargin=1.2cm]
  \item Aplicar la migración V9 para eliminar el campo legado \texttt{libros.autor}
        (String), una vez que el frontend consuma la relación N:M real con la entidad
        Autor.
  \item Implementar rate limiting en los endpoints de autenticación (por ejemplo, con
        Bucket4j y Redis), documentado como gap real en la auditoría OWASP (A04/A07).
  \item Actualizar \texttt{jjwt} de 0.12.6 a 0.13.0 por tocar directamente la superficie
        de autenticación, y planificar la migración a Spring Boot 4 en un bloque dedicado.
  \item Mitigar el riesgo de cache stampede del ADR-003 (por ejemplo, con
        \texttt{@Cacheable(sync=true)} o invalidación selectiva por página).
  \item Agregar a la API REST un endpoint de búsqueda por ISBN equivalente al del servicio
        SOAP, para comparar ambos servicios sobre el mismo contrato de datos.
  \item Completar la PWA del frontend y registrar el reporte Lighthouse como evidencia.
\end{itemize}

\newpage

% ---------------------------------------------------------------------------
% 9. REFERENCIAS
% ---------------------------------------------------------------------------
\section{Referencias}

La guía de práctica exige referencias en norma \textbf{IEEE} y el enunciado del SGA exige
norma \textbf{APA}. Siguiendo la regla del equipo de cumplir ambos requisitos cuando el
SGA añade una exigencia que la guía no contradice (decisión documentada en
\texttt{DECISIONES-GUIA-VS-ENUNCIADO-U4.md}), las citas en el cuerpo se presentan en
autor-año (APA) y, para las fuentes que la guía exige como mínimo, también con su número
de cita IEEE \texttt{[n]} en el mismo lugar del texto; ese número corresponde a la lista
IEEE que aparece al final de esta sección. Se presentan los dos listados: el primero en
APA (para el SGA) y el segundo con las referencias mínimas IEEE que exige la guía.

\subsection*{Referencias en norma APA (exigidas por el SGA)}

\begin{itemize}[leftmargin=1.5cm, itemindent=-1.5cm, itemsep=5pt]
  \item \hypertarget{ref:apache}{Apache Software Foundation. (s.f.). \textit{ab --- Apache HTTP server benchmarking tool}.} Recuperado el 10 de agosto de 2026, de \url{https://httpd.apache.org/docs/current/programs/ab.html}
  \item \hypertarget{ref:breeding}{Breeding, M. (2017). Chapter 2. Koha. The original open source ILS. \textit{Library Technology Reports}, \textit{53}(6).} \url{https://journals.ala.org/index.php/ltr/article/view/6405/8452}
  \item \hypertarget{ref:docker}{Docker. (s.f.). \textit{Docker docs}.} Recuperado el 10 de agosto de 2026, de \url{https://docs.docker.com/}
  \item \hypertarget{ref:fielding2000}{Fielding, R. T. (2000). \textit{Architectural styles and the design of network-based software architectures} [Tesis doctoral, University of California, Irvine].} \url{https://ics.uci.edu/~fielding/pubs/dissertation/rest_arch_style.htm}
  \item \hypertarget{ref:fielding2022}{Fielding, R., Nottingham, M., \& Reschke, J. (Eds.). (2022). \textit{HTTP Semantics} (RFC 9110). Internet Engineering Task Force.} \url{https://www.rfc-editor.org/rfc/rfc9110}
  \item \hypertarget{ref:fowler}{Fowler, M. (2002). \textit{Patterns of Enterprise Application Architecture}.} Addison-Wesley Professional.
  \item \hypertarget{ref:google}{Google. (s.f.). \textit{Learn PWA}.} Recuperado el 10 de agosto de 2026, de \url{https://web.dev/progressive-web-apps/}
  \item \hypertarget{ref:jones2015}{Jones, M., Bradley, J., \& Sakimura, N. (2015). \textit{JSON Web Token (JWT)} (RFC 7519). Internet Engineering Task Force.} \url{https://datatracker.ietf.org/doc/html/rfc7519}
  \item \hypertarget{ref:netlify}{Netlify. (s.f.). \textit{Jamstack}.} Recuperado el 10 de agosto de 2026, de \url{https://jamstack.org/}
  \item \hypertarget{ref:nygard}{Nygard, M. (2018). \textit{Release It! Design and deploy production-ready software} (2.ª ed.). Pragmatic Bookshelf.}
  \item \hypertarget{ref:openapi}{OpenAPI Initiative. (2023). \textit{OpenAPI Specification 3.1}.} \url{https://spec.openapis.org/oas/latest.html}
  \item \hypertarget{ref:otwell}{Otwell, T. (2024). \textit{Laravel 11.x documentation}.} Recuperado el 10 de agosto de 2026, de \url{https://laravel.com/docs/11.x}
  \item \hypertarget{ref:owasp}{OWASP Foundation. (2021). \textit{OWASP Top 10:2021}.} \url{https://owasp.org/Top10/}
  \item \hypertarget{ref:postgres}{PostgreSQL Global Development Group. (s.f.). \textit{PostgreSQL documentation}.} Recuperado el 10 de agosto de 2026, de \url{https://www.postgresql.org/docs/}
  \item \hypertarget{ref:redis}{Redis. (s.f.). \textit{Redis documentation}.} Recuperado el 10 de agosto de 2026, de \url{https://redis.io/docs/latest/}
  \item \hypertarget{ref:richardson}{Richardson, L., \& Ruby, S. (2007). \textit{RESTful Web Services}.} O'Reilly Media.
  \item \hypertarget{ref:springboot}{Spring Boot. (s.f.). \textit{Spring Boot reference documentation}.} Recuperado el 10 de agosto de 2026, de \url{https://docs.spring.io/spring-boot/}
  \item \hypertarget{ref:springsecurity}{Spring Security. (s.f.). \textit{Spring Security reference}.} Recuperado el 10 de agosto de 2026, de \url{https://docs.spring.io/spring-security/}
  \item \hypertarget{ref:washizaki}{Washizaki, H. (Ed.). (2024). \textit{Guide to the Software Engineering Body of Knowledge (SWEBOK), Version 4.0}. IEEE Computer Society.} \url{https://www.computer.org/education/bodies-of-knowledge/software-engineering}
  \item \hypertarget{ref:w3c}{World Wide Web Consortium. (2000). \textit{SOAP Version 1.1} (W3C Note).} \url{https://www.w3.org/TR/soap/}
\end{itemize}

\subsection*{Referencias mínimas en norma IEEE (exigidas por la guía de práctica)}

\begin{enumerate}[label={[\arabic*]}, leftmargin=1.2cm, itemsep=5pt]
\item \hypertarget{ieee1}{T. Otwell, ``Laravel 11.x Documentation,'' Laravel LLC, 2024. [Online]. Available: \url{https://laravel.com/docs/11.x}}
\item \hypertarget{ieee2}{M. Fowler, \textit{Patterns of Enterprise Application Architecture}. Boston, MA, USA: Addison-Wesley, 2002, ISBN: 978-0-321-12521-7.}
\item \hypertarget{ieee3}{R. T. Fielding, ``Architectural Styles and the Design of Network-Based Software Architectures,'' Ph.D. dissertation, Univ. of California, Irvine, CA, USA, 2000. [Online]. Available: \url{https://ics.uci.edu/~fielding/pubs/dissertation/top.htm}}
\item \hypertarget{ieee4}{L. Richardson and S. Ruby, \textit{RESTful Web Services}. Sebastopol, CA, USA: O'Reilly Media, 2007, ISBN: 978-0-596-52926-0.}
\item \hypertarget{ieee5}{H. Washizaki, Ed., \textit{Guide to the Software Engineering Body of Knowledge (SWEBOK), Version 4.0}. IEEE Computer Society, 2024.}
\item \hypertarget{ieee6}{OWASP Foundation, ``OWASP Top 10:2021,'' 2021. [Online]. Available: \url{https://owasp.org/Top10/}}
\item \hypertarget{ieee7}{M. T. Nygard, \textit{Release It! Design and Deploy Production-Ready Software}, 2nd ed. Pragmatic Bookshelf, 2018, ISBN: 978-1-68050-239-8.}
\end{enumerate}

\newpage

% ---------------------------------------------------------------------------
% ANEXO A
% ---------------------------------------------------------------------------
\section*{Anexo A: Estructura del repositorio (versión final en main)}
\addcontentsline{toc}{section}{Anexo A: Estructura del repositorio}

La estructura final del repositorio en \texttt{main}, generada a partir del árbol real del
proyecto (se omiten \texttt{node\_modules}, \texttt{target} y \texttt{.git}), es la
siguiente:

\begin{lstlisting}[basicstyle=\ttfamily\scriptsize, columns=fullflexible, caption={Estructura del repositorio.}]
.
+-- backend/
|   +-- src/main/java/ec/edu/uteq/pfcbackend/
|   |   +-- client/          OpenLibraryClient, OpenLibraryResult
|   |   +-- config/          SecurityConfig, RedisCacheConfig, WebClientConfig,
|   |   |                    WebServiceConfig, OpenApiConfig, CacheablePage
|   |   +-- controller/      Auth, Libro, Autor, Editorial, Idioma, EstadoLibro,
|   |   |                    DocsRedirect
|   |   +-- dto/             ApiResponse, LibroEnriquecidoResponse, LoginRequest, ...
|   |   +-- entity/          Libro, Autor, Editorial, Idioma, EstadoLibro, Usuario
|   |   +-- exception/       GlobalExceptionHandler, BusinessException, ...
|   |   +-- repository/      LibroRepository, UsuarioRepository, ...
|   |   +-- security/        JwtService, JwtAuthenticationFilter, TokenBlacklistService
|   |   +-- service/         LibroService(+Impl), AutorService(+Impl), ...
|   |   +-- ws/              LibroCatalogoEndpoint (SOAP), LibroNoEncontradoSoapException
|   +-- src/main/resources/
|   |   +-- application.yml
|   |   +-- db/migration/    V1..V8 (Flyway)
|   |   +-- libro-catalogo.xsd
|   +-- src/test/java/ec/edu/uteq/pfcbackend/
|   |   +-- integration/     BaseIntegracionTest, AuthIntegracionTest, LibroIntegracionTest
|   |   +-- client/          OpenLibraryClientTest
|   |   +-- repository/      LibroRepositoryTest
|   |   +-- service/         LibroServiceImplTest, AutorServiceImplTest
|   +-- Dockerfile, pom.xml, mvnw
+-- frontend/src/app/
|   +-- core/                guards/auth.guard, interceptors, models, services
|   +-- pages/               login/, libros/, libro-form/, libro-detalle/, autores/, autor-form/
|   +-- shared/               notification-banner
+-- nginx/nginx.conf         reverse proxy de entrada (unico puerto publicado)
+-- docker-compose.yml       5 servicios: postgres, redis, app, frontend, nginx
+-- docs/
|   +-- adr/                 ADR-001 a ADR-005
|   +-- arquitectura/        C4, diagrama de clases y ER (puml + png)
|   +-- informe/             informe-unidad-iv.tex, apache-bench.md, auditoria-owasp.md,
|   |                        reflexion-maria.md, reflexion-marlon.md, reflexion-kevin.md,
|   |                        lighthouse/
|   +-- postman/             SGB-API.postman_collection.json (29 requests)
+-- .env.example             JWT_SECRET, COOKIE_SECURE, etc.
+-- README.md
\end{lstlisting}

\newpage

% ---------------------------------------------------------------------------
% ANEXO C
% ---------------------------------------------------------------------------
\section*{Anexo C: Análisis y discusión}
\addcontentsline{toc}{section}{Anexo C: Análisis y discusión}

Las cuatro preguntas que se responden a continuación son las del Anexo C de la guía de
práctica (\textit{guia\_practica\_experimental\_IV\_app\_web.pdf}). Cada nivel de Bloom
se indica entre corchetes, como en el enunciado. Las respuestas se basan únicamente en
la evidencia de este repositorio (ADRs, benchmark, auditoría OWASP y código real).

\subsection*{Pregunta 1 (Evaluación / Bloom nivel 6)}

\textbf{Revisando el PFC completo (PE-U1 a PE-U4): ¿cuál fue la decisión técnica más
difícil que tomó el equipo? ¿Qué alternativas consideraron? ¿Con el conocimiento actual,
tomarían la misma decisión o una diferente? Justifique con criterios técnicos
referenciados.}

La decisión más difícil fue la elección del framework de backend, porque concentraba los
mayores costos de migración y porque las tres alternativas que el propio PDF declara
válidas (Laravel 11.x, ASP.NET Core 8.x y Spring Boot 3.x) eran técnicamente razonables
(ADR-005). Se consideraron las tres: Laravel se descartó por falta de experiencia real
del equipo con PHP y por el tipado dinámico (un error de tipo en un DTO con reglas de
autorización por rol solo se detectaría en ejecución); ASP.NET Core se descartó pese a ser
la opción de mayor rendimiento bruto en TechEmpower Round 23 (~609\,966 req/s frente a
~243\,639 req/s de Spring en la prueba Fortunes), porque el rendimiento no era la
restricción activa de este proyecto. Se eligió Spring Boot por el dominio previo del
equipo con Java, el tipado estático y un ecosistema de testing maduro (Testcontainers,
JUnit 5). Con el conocimiento actual el equipo tomaría la misma decisión: el resultado
del proyecto (5 recursos CRUD, JWT, cache-aside, SOAP contract-first, 48 tests y un
Docker Compose de 5 servicios) no mostró que el rendimiento de Spring fuera una limitación
real, y la base de conocimiento previa fue determinante con un plazo acotado. La evidencia
de desempeño (P99 de 36 ms con cache frío y 30 ms con cache caliente en la prueba de carga)
respalda retrospectivamente que el factor de elección no debía ser el rendimiento bruto.

\subsection*{Pregunta 2 (Síntesis / Bloom nivel 5)}

\textbf{Los resultados de Apache Bench mostraron \_\_\_ req/s con \_\_\_\% de error. Si se
exigiera que el sistema soporte 500 usuarios concurrentes con tasa de error 0\%, ¿qué
cambios de arquitectura y código serían necesarios? Dibuje la arquitectura escalada.}

Los resultados de Apache Bench mostraron \textbf{802.66 req/s} con la cache fría y
\textbf{986.80 req/s} con la cache caliente, con \textbf{0\% de error} en ambas corridas
(\texttt{n=500}, \texttt{c=20}, endpoint \texttt{GET /api/v1/libros} contra Nginx). Para
soportar 500 usuarios concurrentes con 0\% de error, el cambio de mayor impacto no es de
código sino de topología, porque la autenticación JWT stateless ya permite escalado
horizontal sin sesiones pegajosas. Los cambios serían:

\begin{itemize}[itemsep=2pt]
  \item \textbf{Escalar el plano de aplicación:} varias réplicas de la app Spring Boot
  detrás de Nginx en modo \textit{load balancer} (\texttt{upstream} con \texttt{round-robin}),
  en vez de una sola instancia. El JWT \texttt{HttpOnly} (ADR-003, dual cookie + header
  \texttt{Authorization}) funciona igual en cualquier réplica porque el servidor no guarda
  estado de sesión.
  \item \textbf{Alta disponibilidad de Redis:} hoy la cache-aside y la blacklist de
  \texttt{jti} comparten una sola instancia de Redis, reconocido en el ADR-003 como un
  punto único de fallo. Para 500 concurrentes se requeriría Redis Sentinel o Cluster, y
  separar la blacklist de la cache en instancias distintas.
  \item \textbf{Escalar PostgreSQL:} aumentar el pool de conexiones (\texttt{HikariCP}),
  usar un \textit{connection pooler} intermedio (PgBouncer) y añadir una réplica de
  lectura; las escrituras son infrecuentes frente a las lecturas del catálogo.
  \item \textbf{Mitigar el \textit{cache stampede}:} el ADR-003 documenta que hoy el riesgo
  existe y no está mitigado (TTL que expira con muchas lecturas simultáneas y
  \texttt{allEntries=true} en cada escritura). A 500 usuarios concurrentes se justificaría
  \texttt{@Cacheable(sync=true)} para serializar la reconstrucción de una key.
  \item \textbf{\textit{Rate limiting} hacia la API externa:} el consumo de Open Library
  pasa por el mismo \texttt{WebClient} (sección de desarrollo); a esa concurrencia se
  requeriría un límite de peticiones por ventana de tiempo y un TTL de cache acorde al
  dominio (los 5 minutos actuales del cache-aside están dentro del rango que la guía
  sugiere para datos que cambian poco).
\end{itemize}

Arquitectura escalada propuesta:

\begin{lstlisting}[basicstyle=\ttfamily\footnotesize]
Cliente (SPA Angular 17+)  ->  https://localhost:443 (Nginx LB, round-robin)
                                        |
                  +---------------------+---------------------+
                  |                     |                     |
               app-1                app-2                 app-N   (Spring Boot,
                  |                     |                     |      JWT stateless)
                  +---------------------+---------------------+
                                        |
                        +---------------+---------------+
                        |                               |
                  Redis Cluster                    PostgreSQL 16
             (cache-aside + blacklist)          (pool + replica
                        |                      de lectura + PgBouncer)
                        |
             API externa: Open Library (WebClient, TTL, rate limit)
\end{lstlisting}

\subsection*{Pregunta 3 (Análisis / Bloom nivel 4)}

\textbf{Analice los 6 principios REST de Fielding contra la API REST del PFC: ¿cuáles
cumple completamente, cuáles parcialmente y cuáles no? Para los principios no cumplidos:
¿cuál es el costo de no cumplirlos en este contexto específico?}

\begin{table}[H]
\centering
\caption{Cumplimiento de los 6 principios de Fielding en la API del PFC.}
\label{tab:fielding-annexo}
\footnotesize
\setlength{\tabcolsep}{4pt}
\renewcommand{\arraystretch}{1.3}
\begin{tabularx}{\textwidth}{lXX}
\toprule
Principio & Cumplimiento & Evidencia en el repo \\
\midrule
Client-Server & Completo & Angular 17+ consume la API; el servidor no renderiza la vista \\
Stateless & Completo & JWT firmado en cada petición; blacklist de \texttt{jti} en Redis (ADR-003) \\
Cacheable & Parcial & Cache-aside en Redis (TTL 5 min); sin cabecera \texttt{Cache-Control} hacia el cliente \\
Uniform Interface & Parcial & URIs semánticas (\texttt{/api/v1/libros}) y verbos HTTP; sin HATEOAS/hipermedia \\
Layered System & Completo & Nginx como reverse proxy delante de la app; el cliente solo ve Nginx \\
Code-On-Demand & No & No se envían scripts descargables; es un principio opcional en REST \\
\bottomrule
\end{tabularx}
\end{table}

\textbf{Cacheable (parcial):} las respuestas se cachean en el servidor (Redis), pero la
API no emite la cabecera \texttt{Cache-Control}, por lo que el navegador o un CDN no
pueden reutilizar la respuesta. El costo en este contexto es bajo: el consumidor real es
una SPA del propio equipo y el cache-aside ya reduce la carga de PostgreSQL; un CDN solo
aportaría si hubiera usuarios externos dispersos.

\textbf{Uniform Interface (parcial):} faltan los \textit{links} de hipermedia (HATEOAS);
las respuestas JSON no enlazan los recursos relacionados. El costo es moderado: el cliente
conoce las URIs de forma estática (convención del equipo), por lo que no se pierde
adaptabilidad, pero un consumidor tercero tendría que leer la documentación OpenAPI en vez
de descubrir los recursos en la propia respuesta. Para un PFC con consumidores conocidos
es aceptable.

\textbf{Code-On-Demand (no cumplido):} REST lo define como opcional y el equipo no lo usa
(la SPA transfiere el código en build, no por HTTP dinámico). El costo es nulo: ningún
escenario del PFC requiere ejecutar código provisto por el servidor en tiempo de ejecución.

\subsection*{Pregunta 4 (Evaluación / Bloom nivel 6)}

\textbf{Comparando el PFC con un sistema web de producción real del mismo dominio (buscar
un caso de estudio publicado), ¿qué características de calidad (ISO/IEC 25010:2023) están
cubiertas por el PFC y cuáles requerirían trabajo adicional para ser comercialmente
viables? ¿Cuánto tiempo estimaría para hacerlo comercialmente viable?}

Tomando el catálogo de una biblioteca como dominio, las características de ISO/IEC
25010:2023 que el PFC cubre con evidencia son: \textbf{adecuación funcional} (5 recursos
CRUD, autenticación con roles, SOAP y REST), \textbf{seguridad} parcial (auditoría OWASP
A01--A07 con 6 hallazgos mitigados y 2 brechas documentadas), \textbf{eficiencia de
desempeño} en carga moderada (P99 30--36 ms a 20 concurrentes) y \textbf{mantenibilidad}
(arquitectura en capas, ADRs, 48 tests y 60\,\% de cobertura de instrucciones medida por
JaCoCo en el último build).

Requerirían trabajo adicional para ser comercialmente viable: \textbf{fiabilidad} (sin
monitoreo, \textit{health checks} ni respaldo/restauración de datos documentados),
\textbf{usabilidad} (sin auditoría de accesibilidad WCAG ni pruebas con usuarios reales),
\textbf{portabilidad} parcial (Docker Compose sí, pero sin \textit{secrets} en el
repositorio ni \textit{pipeline} CI/CD automatizado) y \textbf{seguridad} restante (las
dos brechas de la auditoría y \textit{rate limiting} frente a consumo abusivo).

Como caso de estudio publicado del mismo dominio, Koha (\hyperlink{ref:breeding}{Breeding, 2017}) es el
primer sistema integrado de gestión bibliotecaria (ILS) de código abierto, en producción
ininterrumpida desde el año 2000 (creado originalmente para el Horowhenua Library Trust
en Nueva Zelanda) y usado hoy en bibliotecas de todo tamaño, incluyendo consorcios
completos — una escala de producción real de 25 años que este PFC no alcanza ni pretende
alcanzar en un semestre. La comparación en tres características de ISO/IEC 25010:2023
deja ver la brecha real, no solo la esperada:

\begin{itemize}[itemsep=2pt]
  \item \textbf{Adecuación funcional:} Koha cubre módulos de circulación, catalogación
  bajo estándares bibliotecarios reales (MARC 21, UNIMARC, Z39.50/SRU), adquisiciones,
  publicaciones seriadas y un OPAC público, además de protocolos de integración con
  hardware externo (SIP2/NCIP para autopréstamo y RFID). Este PFC cubre 5 recursos CRUD
  genéricos y dos protocolos de exposición (REST y SOAP), sin ningún estándar
  bibliotecario real, porque no era el objetivo del curso.
  \item \textbf{Fiabilidad:} Koha tiene 25 años de operación en producción real, con un
  proceso de release formal respaldado por una comunidad activa. Este PFC no tiene
  monitoreo más allá de los healthchecks de Docker, ni una estrategia de
  respaldo/restauración de PostgreSQL documentada.
  \item \textbf{Portabilidad/Interoperabilidad:} Koha se integra con sistemas de terceros
  mediante protocolos bibliotecarios estándar (SIP2, Z39.50/SRU para catalogación
  compartida entre bibliotecas). Este PFC integra un único servicio externo (Open
  Library, vía REST) y no implementa ningún estándar de interoperabilidad bibliotecaria.
\end{itemize}

Desglose temporal estimado para acercar este PFC a un nivel de madurez comparable, más
granular que el rango de 2 a 3 meses ya mencionado: 1 semana para monitoreo básico
(métricas y logs centralizados), 1 semana para respaldos automáticos de PostgreSQL, 2
semanas para cerrar las 2 brechas de la auditoría OWASP (rate limiting y actualización de
\texttt{jjwt}), 2 semanas para pruebas de accesibilidad WCAG y ajustes de UI, y 2 semanas
para un pipeline CI/CD con despliegue de múltiples réplicas y TLS gestionado (8 semanas
en total, con una persona dedicada). Ninguna de esas 8 semanas cierra la brecha de
interoperabilidad bibliotecaria real (MARC, Z39.50, SIP2) que Koha sí tiene: implementar
esos estándares es un proyecto de dominio en sí mismo, no una tarea de endurecimiento del
sistema actual, y por eso se mantiene fuera de la estimación de viabilidad comercial de
este PFC en su forma actual.

\end{document}
